\section{Symbols and projective points over \texorpdfstring{$p$}{p}th powers}
\label{sec:symbols}

Let $F/\kappa$ be a finitely generated extension of an algebraically closed
field of characteristic $p$, and put
\[
 C=F^p,\qquad d=\trdeg(F/\kappa).
\]
Since $(F^\times)^p=C^\times$, the underlying set of $V_F$ is the point set
of the projective space $\PP_C(F)$: the class $[f]$ represents the
one-dimensional $C$-subspace $Cf$.  These projective points also carry the
$\Fp$-vector-space structure $[f]+[g]=[fg]$, whose zero is $[1]$.
We will use both structures, keeping their two scalar fields distinct.
Our first task is to extract from a vanishing symbol the field-theoretic
constraint that will allow us to recover projective lines.

\begin{lemma}\label{lem:p-degree}
With the notation above, $[F:C]=p^d$.
\end{lemma}

\begin{proof}
Choose a transcendence basis $t_1,\ldots,t_d$ over $\kappa$, and write
$F_0=\kappa(t_1,\ldots,t_d)$.  Perfectness of $\kappa$ gives
$F_0^p=\kappa(t_1^p,\ldots,t_d^p)$.  The monomials
$t_1^{i_1}\cdots t_d^{i_d}$ with $0\leq i_j<p$ form a basis of $F_0$ over
$F_0^p$: group polynomial exponents modulo $p$, and for rational functions
use $a/b=ab^{p-1}/b^p$.  Linear independence follows by clearing
denominators and separating the exponent classes.
Thus $[F_0:F_0^p]=p^d$.

The extension $F/F_0$ is finite, and Frobenius induces an isomorphism of
extensions $F/F_0\simeq F^p/F_0^p$.  Computing $[F:F_0^p]$ in the two
towers through $F_0$ and $F^p$ gives
\[
 [F:F_0]p^d=[F:F^p][F^p:F_0^p]
           =[F:F^p][F:F_0].
\]
Cancellation proves the assertion without a separability assumption.
\end{proof}

A list $s_1,\ldots,s_r\in F$ is \emph{$p$-independent over $C$} if
$[C(s_1,\ldots,s_r):C]=p^r$.  A $p$-independent list generating $F$ over
$C$ is called a \emph{$p$-basis}.  Its monomials supply derivations that
record one exponent at a time.

\begin{lemma}\label{lem:p-basis}
Let $F$ be any field of characteristic $p$, and put $C=F^p$.
If $C\subseteq M\subseteq F$ and $s\in F\setminus M$, then
$[M(s):M]=p$.  If $[F:C]$ is finite, every $p$-independent list extends to a $p$-basis
$s_1,\ldots,s_e$ of $F/C$, with $[F:C]=p^e$ and monomial basis
\[
 \{s_1^{i_1}\cdots s_e^{i_e}:0\leq i_j<p\}.
\]
For such a basis there are commuting $C$-linear derivations
$D_1,\ldots,D_e$ on $F$ satisfying
\[
 D_j(s_1^{i_1}\cdots s_e^{i_e})
       =i_j s_1^{i_1}\cdots s_e^{i_e}.
\]
On $C(s_j)$ the kernel of $D_j$ is $C$, and its nonzero eigenvectors with
eigenvalue in $C$ are precisely the elements $c s_j^i$, where
$c\in C^\times$ and $0\leq i<p$; their eigenvalues are $i\in\Fp$.
Moreover, $C(s_j)\cap C(s_k)=C$ when $j\ne k$.
\end{lemma}

\begin{proof}
Since $s^p\in C$, the minimal polynomial of $s$ over $M$ divides
$T^p-s^p=(T-s)^p$.  If its degree were $1\leq e<p$, it would be
$(T-s)^e$, and its coefficient $-es$ of $T^{e-1}$ would force $s\in M$.
Thus every strict adjunction has degree $p$.  When $[F:C]$ is finite,
successive adjunctions must terminate, so they extend any
$p$-independent list to a generating list.  If its length is $e$, the
tower formula gives $[F:C]=p^e$, and the usual basis of each simple
extension gives the asserted monomial basis.

Define the $D_j$ on these monomials and extend $C$-linearly.  When two
monomials are multiplied, reducing an exponent by $p$ only removes a
factor in $C$ and leaves its value in $\Fp$ unchanged.  This proves the
product rule.  The operators commute because they are diagonal on the
same basis.  On $C(s_j)$ their eigenvector assertions follow from the
distinct eigenvalues $0,1,\ldots,p-1$ on $1,s_j,\ldots,s_j^{p-1}$.
Finally, the expansions of an element in both $C(s_j)$ and $C(s_k)$
can have only their constant term in common.
\end{proof}

For a $p$-independent pair $s,t$, write $D_s,D_t$ for the corresponding
derivations obtained from any extending $p$-basis.  The following proof
uses a coordinate form of the logarithmic differential symbol;
the much stronger injectivity theorem for that symbol is the
Bloch--Gabber--Kato theorem \citep[Theorem~(2.1)]{BlochKato1986}.
Here we check the needed implication directly.

\begin{lemma}\label{lem:zero-symbol}
For $s,t\in F^\times$, one has
\[
 s\notin C,\qquad \{s,t\}_F=0
 \quad\Longrightarrow\quad t\in C(s).
\]
\end{lemma}

\begin{proof}
Suppose instead that $s\notin C$ and $t\notin C(s)$.  By
Lemma~\ref{lem:p-basis}, the pair $s,t$ is $p$-independent.  Consider
\[
 \beta([f],[g])=
 \frac{D_s f}{f}\frac{D_t g}{g}
 -\frac{D_t f}{f}\frac{D_s g}{g},
 \qquad f,g\in F^\times.
\]
Each logarithmic derivative is additive on products and vanishes on
$C^\times$.  Hence $\beta$ is a well-defined $\Fp$-bilinear map
$V_F\times V_F\to F$.  For $f\ne0,1$, the identities
$D_s(1-f)=-D_s f$ and $D_t(1-f)=-D_t f$ give
$\beta([f],[1-f])=0$.  Thus $\beta$ descends through the Steinberg
relations to an $\Fp$-linear map $W_F\to F$.
But $D_s s=s$, $D_t t=t$, and $D_s t=D_t s=0$, so this map sends
$\{s,t\}_F$ to $1$.  The symbol cannot vanish.
\end{proof}

The symbols arising from an affine line do vanish, by the defining
Steinberg relation itself.

\begin{lemma}\label{lem:affine-symbol}
If $c,d\in C^\times$ and $t,c+dt\in F^\times$, then
\[
 \{t,c+dt\}_F=0.
\]
\end{lemma}

\begin{proof}
Put $q=-dt/c$.  The hypotheses give $q\ne0,1$, while
$[q]=[t]$ and $[1-q]=[c+dt]$, since the omitted factors lie in
$C^\times$.  Therefore the identity is the Steinberg relation
$\{q,1-q\}_F=0$.
\end{proof}
